\documentclass[11pt,a4paper]{article}
\usepackage[utf8]{inputenc}
\usepackage[T1]{fontenc}
\usepackage[french]{babel}
\usepackage{amsmath,amssymb,amsthm}
\usepackage{geometry}
\geometry{hmargin=2.5cm,vmargin=3cm}
\usepackage{tikz}
\usetikzlibrary{cd,positioning,shapes}
\usepackage{listings}
\usepackage{xcolor}

\lstset{
    language=Python,
    basicstyle=\ttfamily\small,
    keywordstyle=\color{blue}\bfseries,
    stringstyle=\color{red},
    commentstyle=\color{gray}\itshape,
    numbers=left,
    numberstyle=\tiny\color{gray},
    breaklines=true,
    frame=single,
    showstringspaces=false,
    literate={é}{{\'e}}1 {è}{{\`e}}1 {à}{{\`a}}1 {ê}{{\^e}}1 {î}{{\^i}}1 {ç}{{\c{c}}}1 {ï}{{\"i}}1 {É}{{\'E}}1
}

\title{\Huge JALON 70 : ESPACES MESURÉS PRODUITS \\ \large Cours magistral, exercices corrigés et travaux pratiques}
\author{\Large Charles EDOU NZE}
\date{\today}

\begin{document}
\maketitle
\tableofcontents
\newpage

---
uuid: "jalon-70"
title: "Espaces mesurés produits"
year: 2
trimester: 6
tags:
  - math/mesure
  - ia/probabilites
prev: "[[Jalon 69 (Démonstration complète du théorème de convergence dominée de Lebesgue.).md]]"
next: "[[Jalon 71 (Théorèmes de Fubini-Tonelli).md]]"
---

\part{Cours Magistral}

\section{Genèse et Intuition}

 Imaginez que vous ayez deux règles. L'une mesure des longueurs sur un axe horizontal ($X$), l'autre mesure des longueurs sur un axe vertical ($Y$). Si vous croisez ces deux règles, vous créez un monde en 2D (le produit $X \times Y$).
    - Une \textbf{Tribu produit}, c'est décider que les "rectangles" (un ensemble de $X$ croisé avec un ensemble de $Y$) sont nos nouvelles briques de base pour mesurer des surfaces.
    - Une \textbf{Mesure produit}, c'est dire que la surface d'un rectangle est simplement sa largeur multipliée par sa hauteur.
    C'est la manière naturelle de construire des mesures complexes à partir de mesures simples.
 La plupart des phénomènes réels dépendent de plusieurs facteurs. Pour calculer le volume d'un objet, ou la probabilité que deux événements indépendants arrivent en même temps, on a besoin de savoir comment "multiplier" les mesures entre elles de manière rigoureuse.
 Un quadrillage. On définit la mesure sur chaque petit carreau, puis on étend cette définition à toutes les formes bizarres que l'on peut construire en assemblant des carreaux.

\section{Définitions, Théorèmes \& Exemples Concrets Immédiats}

Soient $(X_1, \mathcal{F}_1, \mu_1)$ et $(X_2, \mathcal{F}_2, \mu_2)$ deux espaces mesurés.

\subsection{Construction de la Tribu Produit}

> \textbf{Exemple Concret Immédiat :}
> Si $X_1 = X_2 = \mathbb{R}$ munis de la tribu borélienne $\mathcal{B}(\mathbb{R})$, la tribu produit $\mathcal{B}(\mathbb{R}) \otimes \mathcal{B}(\mathbb{R})$ est exactement la tribu borélienne de $\mathbb{R}^2$, notée $\mathcal{B}(\mathbb{R}^2)$.
> Le rectangle géométrique $[0, 1] \times [0, 1]$ est un rectangle mesurable.


> \textbf{Définition 1 (Rectangle mesurable) :}
> On appelle \textbf{rectangle mesurable} toute partie de $X_1 \times X_2$ de la forme $A_1 \times A_2$ où $A_1 \in \mathcal{F}_1$ et $A_2 \in \mathcal{F}_2$.

> \textbf{Définition 2 (Tribu Produit) :}
> La \textbf{tribu produit}, notée $\mathcal{F}_1 \otimes \mathcal{F}_2$, est la tribu engendrée par l'ensemble des rectangles mesurables sur $X_1 \times X_2$.

\subsection{Existence et Unicité de la Mesure Produit}

> \textbf{Exemple Concret Immédiat :}
> Prenons $X_1 = X_2 = \mathbb{R}$ avec $\mu_1 = \mu_2 = \lambda$ (mesure de Lebesgue). Les intervalles $A_1 = [0, 2]$ et $A_2 = [1, 4]$ sont mesurables.
> La mesure produit $\lambda \otimes \lambda$ du rectangle $R = A_1 \times A_2$ est simplement l'aire géométrique :
> $\pi(R) = \lambda([0, 2]) \cdot \lambda([1, 4]) = 2 \cdot 3 = 6$.
> La condition de $\sigma$-finitude est bien vérifiée car $\mathbb{R} = \bigcup_{n \in \mathbb{N}} [-n, n]$.


> \textbf{Théorème (Existence et Unicité) :}
> Si les mesures $\mu_1$ et $\mu_2$ sont \textbf{$\sigma$-finies}, alors il existe une unique mesure $\pi$ sur $(X_1 \times X_2, \mathcal{F}_1 \otimes \mathcal{F}_2)$, notée $\mu_1 \otimes \mu_2$, telle que pour tout rectangle mesurable $A_1 \times A_2$ :
> $$\pi(A_1 \times A_2) = \mu_1(A_1) \cdot \mu_2(A_2)$$

\subsection{Sections Mesurables}

Soit $E \in \mathcal{F}_1 \otimes \mathcal{F}_2$. Pour tout $x \in X_1$, on définit la \textbf{section} de $E$ en $x$ par :
$$E_x = \{ y \in X_2 \mid (x, y) \in E \}$$
*Propriété :* Pour tout $x$, $E_x \in \mathcal{F}_2$ (la section d'un mesurable est mesurable).

\section{Démonstrations}

\subsection{Démonstration : Calcul de la mesure par intégration des sections}

Montrons que $\pi(E) = \int_{X_1} \mu_2(E_x) d\mu_1(x)$.

1. \textbf{Cas des rectangles :} Soit $E = A_1 \times A_2$.
   - Si $x \in A_1$, alors $E_x = A_2$. Sa mesure est $\mu_2(A_2)$.
   - Si $x \notin A_1$, alors $E_x = \emptyset$. Sa mesure est 0.
   - Donc $\mu_2(E_x) = \mu_2(A_2) \mathbf{1}_{A_1}(x)$.
   - En intégrant par rapport à $\mu_1$ : $\int \mu_2(A_2) \mathbf{1}_{A_1}(x) d\mu_1 = \mu_2(A_2) \mu_1(A_1) = \pi(E)$.
2. \textbf{Généralisation par classes monotones :}
   Soit $\mathcal{M}$ l'ensemble des parties $E \in \mathcal{F}_1 \otimes \mathcal{F}_2$ pour lesquelles l'égalité $\pi(E) = \int_{X_1} \mu_2(E_x) d\mu_1(x)$ est vraie.
   - Nous venons de montrer que $\mathcal{M}$ contient tous les rectangles mesurables. L'ensemble des rectangles forme un $\pi$-système (stable par intersection finie).
   - $\mathcal{M}$ contient l'espace entier $X_1 \times X_2$.
   - $\mathcal{M}$ est stable par union dénombrable croissante (par le théorème de convergence monotone appliqué à la suite d'indicatrices correspondantes).
   - $\mathcal{M}$ est stable par différence (si $A \subset B$ et $A, B \in \mathcal{M}$, alors par linéarité de l'intégrale, $B \setminus A \in \mathcal{M}$).
   - D'après le théorème des classes monotones, $\mathcal{M}$ contient la tribu engendrée par le $\pi$-système des rectangles, c'est-à-dire $\mathcal{F}_1 \otimes \mathcal{F}_2$.
3. \textbf{Conclusion :} La formule est vraie pour tout ensemble de la tribu produit. Cela montre que la mesure produit est "cohérente" avec l'intégration couche par couche.





\section{Application en Intelligence Artificielle}

- \textbf{Le Pont Théorique :} En IA, on travaille presque exclusivement dans des espaces produits. Une image de $224 \times 224$ pixels est un point dans un espace produit de dimension 50 176.
- \textbf{Example Concret :}
    - \textbf{Indépendance Statistique :} Deux variables aléatoires $X$ et $Y$ sont indépendantes si et seulement si leur mesure de probabilité jointe $P_{X,Y}$ est égale à la mesure produit des marginales $P_X \otimes P_Y$. C'est la définition mathématique de l'absence de corrélation.
    - \textbf{Modèles Génératifs (GANs, VAEs) :} On essaie souvent de factoriser la distribution des données dans un espace latent où les dimensions sont indépendantes. Topologiquement, cela revient à forcer la mesure latente à être une mesure produit.
    - \textbf{Calcul de l'erreur moyenne (Risk) :} Le risque est une intégrale sur l'espace produit "Données $\times$ Étiquettes". Pour l'estimer, on suppose que les exemples sont "IID" (indépendants et identiquement distribués), ce qui permet d'utiliser la structure de mesure produit.

\section{Liens Sémantiques}

- \textbf{Concepts Précédents requis :} [[Jalon 63 (Définition axiomatique d'une mesure).md]], [[Jalon 62 (Algèbres).md]]
- \textbf{Concepts Futurs dépendants :} [[Jalon 71 (Théorèmes de Fubini-Tonelli).md]], [[Jalon 88 (Indépendance d'événements).md]]
\part{Exercices d'Application et de Concours}
\subsection*{Exercice 1 : Calcul de mesure produit sur des intervalles finis \quad $\bigstar\star\star\star\star$}

\textbf{Énoncé :}
Soit $(\mathbb{R}^2, \mathcal{B}(\mathbb{R}^2), \lambda_2)$ où $\lambda_2 = \lambda \otimes \lambda$ est la mesure de Lebesgue sur $\mathbb{R}^2$.
Calculer $\lambda_2(A)$ pour $A = [1, 3] \times [2, 5]$.

\textbf{Correction :}
Par définition de la mesure produit sur un rectangle mesurable :
1. Les intervalles $A_1 = [1, 3]$ et $A_2 = [2, 5]$ sont des boréliens de $\mathbb{R}$.
2. On applique la formule : $\lambda_2(A_1 \times A_2) = \lambda(A_1) \cdot \lambda(A_2)$.
3. Calcul des mesures 1D : $\lambda([1, 3]) = 3 - 1 = 2$ et $\lambda([2, 5]) = 5 - 2 = 3$.
4. Conclusion : $\lambda_2(A) = 2 \cdot 3 = 6$.

\subsection*{Exercice 2 : Produit de mesures discrètes \quad $\bigstar\bigstar\star\star\star$}

\textbf{Énoncé :}
Soient $X_1 = X_2 = \mathbb{N}$ munis de la tribu $\mathcal{P}(\mathbb{N})$ et de la mesure de comptage $c$.
Montrer que $c \otimes c$ est la mesure de comptage sur $\mathbb{N}^2$. Calculer la mesure du sous-ensemble $D = \{ (n, n) \mid n \in \{1, 2, 3\} \}$.

\textbf{Correction :}
1. Tout sous-ensemble de $\mathbb{N}^2$ peut s'écrire comme une union au plus dénombrable de singletons.
2. Pour un singleton $E = \{(a, b)\} = \{a\} \times \{b\}$, qui est un rectangle mesurable, on a :
   $(c \otimes c)(E) = c(\{a\}) \cdot c(\{b\}) = 1 \cdot 1 = 1$.
3. La mesure $c \otimes c$ attribue un poids de 1 à chaque point de $\mathbb{N}^2$, c'est donc bien la mesure de comptage sur cet espace.
4. L'ensemble $D = \{(1,1), (2,2), (3,3)\}$ est l'union disjointe de 3 singletons.
5. $(c \otimes c)(D) = 1 + 1 + 1 = 3$.

\subsection*{Exercice 3 : La section d'un mesurable \quad $\bigstar\bigstar\bigstar\star\star$}

\textbf{Énoncé :}
Soit $E = \{ (x,y) \in [0,1]^2 \mid x^2 + y^2 \le 1 \}$.
Montrer que $E \in \mathcal{B}(\mathbb{R}) \otimes \mathcal{B}(\mathbb{R})$ et décrire formellement la section $E_x$ pour $x \in \mathbb{R}$.

\textbf{Correction :}
1. La fonction $f(x,y) = x^2 + y^2$ est continue sur $\mathbb{R}^2$, donc mesurable par rapport à $\mathcal{B}(\mathbb{R}^2)$.
2. L'ensemble $E = f^{-1}([0, 1]) \cap ([0,1] \times [0,1])$ est l'intersection de deux fermés, donc un borélien de $\mathbb{R}^2$. On rappelle que $\mathcal{B}(\mathbb{R}^2) = \mathcal{B}(\mathbb{R}) \otimes \mathcal{B}(\mathbb{R})$.
3. Pour un $x \in \mathbb{R}$ fixé :
   - Si $x \notin [0, 1]$, alors $E_x = \emptyset$.
   - Si $x \in [0, 1]$, alors on cherche les $y \in [0, 1]$ tels que $y^2 \le 1 - x^2$. Comme $y \ge 0$, on obtient $0 \le y \le \sqrt{1 - x^2}$.
   - Donc $E_x = [0, \sqrt{1 - x^2}]$.

\subsection*{Exercice 4 : Mesure d'un triangle \quad $\bigstar\bigstar\bigstar\star\star$}

\textbf{Énoncé :}
Calculer la mesure de Lebesgue $\lambda_2$ du triangle $T = \{ (x, y) \in \mathbb{R}^2 \mid x \ge 0, y \ge 0, x+y \le 2 \}$ en intégrant ses sections.

\textbf{Correction :}
1. Les sections $T_x$ sont données par :
   - Si $x < 0$ ou $x > 2$, $T_x = \emptyset$ et $\lambda(T_x) = 0$.
   - Si $x \in [0, 2]$, $T_x = [0, 2-x]$ et $\lambda(T_x) = 2-x$.
2. Par le théorème d'intégration des sections (qui découle de la définition de la mesure produit) :
   $\lambda_2(T) = \int_{\mathbb{R}} \lambda(T_x) d\lambda(x) = \int_0^2 (2-x) dx$.
3. Calcul de l'intégrale :
   $\int_0^2 (2-x) dx = \left[ 2x - \frac{x^2}{2} \right]_0^2 = (4 - 2) - 0 = 2$.

\subsection*{Exercice 5 : Indépendance et Produit de probabilités \quad $\bigstar\bigstar\bigstar\star\star$}

\textbf{Énoncé :}
Soient $X$ et $Y$ deux variables aléatoires réelles. Leur loi jointe $P_{(X,Y)}$ est une probabilité sur $(\mathbb{R}^2, \mathcal{B}(\mathbb{R}^2))$.
Montrer que si $P_{(X,Y)} = P_X \otimes P_Y$, alors $P(X \in A, Y \in B) = P(X \in A)P(Y \in B)$ pour tous boréliens $A$ et $B$.

\textbf{Correction :}
1. L'événement $\{X \in A, Y \in B\}$ correspond formellement à $(X, Y) \in A \times B$.
2. La probabilité de cet événement est $P_{(X,Y)}(A \times B)$.
3. Par hypothèse, $P_{(X,Y)}$ est la mesure produit de $P_X$ et $P_Y$.
4. Par définition de la mesure produit sur un rectangle mesurable :
   $P_X \otimes P_Y(A \times B) = P_X(A) \cdot P_Y(B)$.
5. Or, par définition des lois marginales, $P_X(A) = P(X \in A)$ et $P_Y(B) = P(Y \in B)$.
6. D'où le résultat : $P(X \in A, Y \in B) = P(X \in A)P(Y \in B)$.

\subsection*{Exercice 6 : Produit d'une mesure discrète et continue \quad $\bigstar\bigstar\bigstar\bigstar\star$}

\textbf{Énoncé :}
Soit $c$ la mesure de comptage sur $(\mathbb{R}, \mathcal{P}(\mathbb{R}))$ et $\lambda$ la mesure de Lebesgue sur $(\mathbb{R}, \mathcal{B}(\mathbb{R}))$. On munit $\mathbb{R}^2$ de la tribu $\mathcal{P}(\mathbb{R}) \otimes \mathcal{B}(\mathbb{R})$.
Soit la diagonale $\Delta = \{(x,x) \mid x \in [0,1]\}$. Calculer les deux intégrales itérées $\int (\int \mathbf{1}_\Delta(x,y) d\lambda(y)) dc(x)$ et $\int (\int \mathbf{1}_\Delta(x,y) dc(x)) d\lambda(y)$.
Pourquoi sont-elles différentes ?

\textbf{Correction :}
1. Première intégrale : On intègre d'abord sur $y$. Pour un $x \in [0,1]$ fixé, la section $\Delta_x$ est le singleton $\{x\}$. Sa mesure de Lebesgue est $\lambda(\{x\}) = 0$.
   Donc $\int \mathbf{1}_\Delta(x,y) d\lambda(y) = 0$. L'intégrale extérieure par rapport à $c$ donne $0$.
2. Seconde intégrale : On intègre d'abord sur $x$. Pour un $y \in [0,1]$ fixé, la section "verticale" est le singleton $\{y\}$. Sa mesure de comptage est $c(\{y\}) = 1$.
   Donc $\int \mathbf{1}_\Delta(x,y) dc(x) = 1$ pour $y \in [0,1]$, et $0$ sinon. L'intégrale extérieure est $\int_{[0,1]} 1 d\lambda(y) = 1$.
3. Les deux intégrales donnent des résultats différents ($0 \neq 1$) car le théorème de Fubini ne s'applique pas. L'espace mesuré $(\mathbb{R}, \mathcal{P}(\mathbb{R}), c)$ n'est pas $\sigma$-fini.

\subsection*{Exercice 7 : Génération de la tribu produit \quad $\bigstar\bigstar\bigstar\bigstar\star$}

\textbf{Énoncé :}
Montrer que si $\mathcal{C}_1$ engendre $\mathcal{F}_1$ et $\mathcal{C}_2$ engendre $\mathcal{F}_2$, et s'il existe des suites d'ensembles de $\mathcal{C}_i$ qui recouvrent $X_i$, alors l'ensemble des rectangles $C_1 \times C_2$ (avec $C_i \in \mathcal{C}_i$) engendre la tribu produit $\mathcal{F}_1 \otimes \mathcal{F}_2$.

\textbf{Correction :}
1. Soit $\mathcal{G}$ la tribu engendrée par les rectangles $C_1 \times C_2$. Il est clair que $\mathcal{G} \subset \mathcal{F}_1 \otimes \mathcal{F}_2$ car $C_1 \times C_2 \in \mathcal{F}_1 \otimes \mathcal{F}_2$.
2. Montrons l'inclusion inverse. Pour un $C_2 \in \mathcal{C}_2$ fixé, considérons $\mathcal{H}_1 = \{ A \in \mathcal{F}_1 \mid A \times C_2 \in \mathcal{G} \}$.
3. $\mathcal{H}_1$ contient $\mathcal{C}_1$ par définition. De plus, on vérifie facilement que $\mathcal{H}_1$ est une tribu (passage au complémentaire, union dénombrable).
4. Donc $\mathcal{H}_1 = \mathcal{F}_1$. Ainsi, pour tout $A \in \mathcal{F}_1$ et tout $C_2 \in \mathcal{C}_2$, $A \times C_2 \in \mathcal{G}$.
5. On répète le même argument. Pour un $A \in \mathcal{F}_1$ fixé, soit $\mathcal{H}_2 = \{ B \in \mathcal{F}_2 \mid A \times B \in \mathcal{G} \}$.
6. $\mathcal{H}_2$ contient $\mathcal{C}_2$ et est une tribu, donc $\mathcal{H}_2 = \mathcal{F}_2$.
7. Ainsi, tout rectangle $A \times B \in \mathcal{F}_1 \times \mathcal{F}_2$ appartient à $\mathcal{G}$, ce qui prouve que $\mathcal{F}_1 \otimes \mathcal{F}_2 = \mathcal{G}$.

\subsection*{Exercice 8 : Le problème du produit infini non mesurable \quad $\bigstar\bigstar\bigstar\bigstar\star$}

\textbf{Énoncé :}
Montrer par un contre-exemple simple que l'union de deux rectangles mesurables n'est pas nécessairement un rectangle mesurable.

\textbf{Correction :}
1. Soit $X_1 = X_2 = \mathbb{R}$.
2. Considérons $R_1 = [0, 1] \times [0, 1]$ et $R_2 = [2, 3] \times [2, 3]$.
3. Supposons par l'absurde que $R_1 \cup R_2 = A \times B$.
4. Le point $(0, 0)$ appartient à $R_1$, donc à $A \times B$. Ainsi $0 \in A$ et $0 \in B$.
5. Le point $(2, 2)$ appartient à $R_2$, donc à $A \times B$. Ainsi $2 \in A$ et $2 \in B$.
6. Si $A \times B$ est un produit cartésien contenant $0$ et $2$ en $X$, et $0$ et $2$ en $Y$, il doit contenir le point $(0, 2)$.
7. Or $(0, 2)$ n'appartient ni à $R_1$ ni à $R_2$. C'est une contradiction.
8. L'ensemble des rectangles mesurables n'est pas stable par union, c'est pourquoi on doit considérer la tribu *engendrée* par ces rectangles.

\subsection*{Exercice 9 : Densité d'une mesure produit (cas à densité) \quad $\bigstar\bigstar\bigstar\bigstar\bigstar$}

\textbf{Énoncé :}
Si $\mu_1$ a pour densité $f_1$ par rapport à $\lambda$ et $\mu_2$ a pour densité $f_2$ par rapport à $\lambda$, montrer que la mesure produit $\mu_1 \otimes \mu_2$ a pour densité la fonction $(x,y) \mapsto f_1(x)f_2(y)$ par rapport à $\lambda \otimes \lambda$.

\textbf{Correction :}
1. Soit un rectangle mesurable $A \times B$.
2. Par définition, $(\mu_1 \otimes \mu_2)(A \times B) = \mu_1(A) \cdot \mu_2(B)$.
3. Comme les mesures ont des densités : $\mu_1(A) = \int_A f_1(x) dx$ et $\mu_2(B) = \int_B f_2(y) dy$.
4. Donc $(\mu_1 \otimes \mu_2)(A \times B) = \left( \int_A f_1(x) dx \right) \left( \int_B f_2(y) dy \right)$.
5. Les fonctions étant positives, et les variables séparables, cela correspond exactement à l'intégrale double sur $A \times B$ :
   $\int_{A \times B} f_1(x)f_2(y) d(\lambda \otimes \lambda)(x,y)$.
6. Cette égalité étant vraie sur tous les rectangles (qui forment un $\pi$-système engendrant la tribu), elle s'étend par le théorème d'extension de Carathéodory (ou lemme de classe monotone) à tout borélien de $\mathbb{R}^2$.
7. La densité est donc bien $f(x,y) = f_1(x)f_2(y)$.

\subsection*{Exercice 10 : Mesure de Dirac produit \quad $\bigstar\bigstar\bigstar\bigstar\bigstar$}

\textbf{Énoncé :}
Soit $a \in X_1$ et $b \in X_2$. Soit $\delta_a$ la mesure de Dirac en $a$ sur $(X_1, \mathcal{F}_1)$ et $\delta_b$ la mesure de Dirac en $b$ sur $(X_2, \mathcal{F}_2)$.
Montrer que $\delta_a \otimes \delta_b = \delta_{(a,b)}$ où $\delta_{(a,b)}$ est la mesure de Dirac au point $(a,b)$ sur l'espace produit.

\textbf{Correction :}
1. Soit un rectangle mesurable $A_1 \times A_2$.
2. Par définition de la mesure produit :
   $(\delta_a \otimes \delta_b)(A_1 \times A_2) = \delta_a(A_1) \cdot \delta_b(A_2)$.
3. Calculons le membre de droite :
   - Si $a \in A_1$ et $b \in A_2$, alors $\delta_a(A_1) = 1$ et $\delta_b(A_2) = 1$. Le produit vaut 1.
   - Sinon, l'un des deux (ou les deux) vaut 0. Le produit vaut 0.
4. Or, $a \in A_1$ et $b \in A_2$ est exactement équivalent à $(a,b) \in A_1 \times A_2$.
5. Donc la mesure du rectangle vaut 1 si $(a,b)$ appartient au rectangle, et 0 sinon.
6. C'est exactement la définition de la mesure de Dirac en $(a,b)$, notée $\delta_{(a,b)}$.
7. Par unicité de l'extension de la mesure produit (les mesures étant trivialement $\sigma$-finies car finies), l'égalité $\delta_a \otimes \delta_b = \delta_{(a,b)}$ est vraie sur toute la tribu produit.

\part{Travaux Pratiques et Simulations Algorithmiques}
\subsection*{TP 1 : Approximation numerique d'une integrale 2D par somme de Riemann}
\begin{lstlisting}
# TP 1 : Approximation d'une integrale double

def integrale_2d(f, x_min, x_max, y_min, y_max, nx=100, ny=100):
    dx = (x_max - x_min) / nx
    dy = (y_max - y_min) / ny
    somme = 0.0
    for i in range(nx):
        x = x_min + i * dx + dx/2
        for j in range(ny):
            y = y_min + j * dy + dy/2
            somme += f(x, y) * dx * dy
    return somme

# Test : volume d'une pyramide simple f(x,y) = x*y sur [0,1]x[0,1]
def f(x, y):
    return x * y

valeur = integrale_2d(f, 0, 1, 0, 1)
theorique = 0.25 # (1/2) * (1/2)
print(f"Valeur numerique : {valeur}")
assert abs(valeur - theorique) < 1e-3, "Erreur dans le calcul de l'integrale 2D"
\end{lstlisting}

\subsection*{TP 2 : Mesure d'un domaine par methode de Monte-Carlo}
\begin{lstlisting}
# TP 2 : Mesure de Monte-Carlo

import random

def mesure_domaine(condition, x_min, x_max, y_min, y_max, n_points=100000):
    aire_rectangle = (x_max - x_min) * (y_max - y_min)
    points_dedans = 0

    for _ in range(n_points):
        x = random.uniform(x_min, x_max)
        y = random.uniform(y_min, y_max)
        if condition(x, y):
            points_dedans += 1

    mesure_estimee = aire_rectangle * (points_dedans / n_points)
    return mesure_estimee

# Test : Aire du quart de disque unite
def quart_disque(x, y):
    return x*x + y*y <= 1.0

aire_est = mesure_domaine(quart_disque, 0, 1, 0, 1)
pi_sur_4 = 3.14159265 / 4.0
print(f"Aire estimee : {aire_est}, Theorique : {pi_sur_4}")
assert abs(aire_est - pi_sur_4) < 1e-2, "L'estimation Monte-Carlo est trop loin"
\end{lstlisting}

\subsection*{TP 3 : Verification de l'independance statistique (Produit de probabilites)}
\begin{lstlisting}
# TP 3 : Independance et probabilite jointe

def verifier_independance(lois_jointes, p_x, p_y, epsilon=1e-6):
    for x in lois_jointes:
        for y in lois_jointes[x]:
            p_jointe = lois_jointes[x][y]
            p_produit = p_x[x] * p_y[y]
            if abs(p_jointe - p_produit) > epsilon:
                return False
    return True

# Definition des probabilites
p_x = { 'pile': 0.5, 'face': 0.5 }
p_y = { 1: 1/6, 2: 1/6, 3: 1/6, 4: 1/6, 5: 1/6, 6: 1/6 }

# Construction d'une loi jointe independante
lois_jointes = {}
for x in p_x:
    lois_jointes[x] = {}
    for y in p_y:
        lois_jointes[x][y] = p_x[x] * p_y[y]

indep = verifier_independance(lois_jointes, p_x, p_y)
print(f"Les variables sont independantes : {indep}")
assert indep == True, "La fonction devrait renvoyer True pour un produit parfait"
\end{lstlisting}

\subsection*{TP 4 : Section d'un ensemble 2D}
\begin{lstlisting}
# TP 4 : Construction de sections

def generer_section(ensemble_points, x_cible, tolerance=1e-3):
    section = []
    for (x, y) in ensemble_points:
        if abs(x - x_cible) < tolerance:
            section.append(y)
    return section

# Construction d'un disque discret
points = []
for i in range(-50, 51):
    for j in range(-50, 51):
        x = i / 50.0
        y = j / 50.0
        if x*x + y*y <= 1.0:
            points.append((x, y))

# Extraction de la section en x=0 (devrait aller de y=-1 a y=1)
sec = generer_section(points, 0.0)
print(f"Taille de la section en 0 : {len(sec)}")
assert len(sec) == 101, f"Taille inattendue {len(sec)}"
assert min(sec) == -1.0 and max(sec) == 1.0, "La section ne couvre pas [-1, 1]"
\end{lstlisting}

\subsection*{TP 5 : Calcul de la densite jointe de deux variables independantes}
\begin{lstlisting}
# TP 5 : Calcul de densite produit

import math

def densite_normale(x, mu=0, sigma=1):
    return (1.0 / (sigma * math.sqrt(2 * math.pi))) * math.exp(-0.5 * ((x - mu)/sigma)**2)

def densite_jointe_independante(x, y, mu_x, sig_x, mu_y, sig_y):
    # La densite jointe est le produit des densites marginales
    return densite_normale(x, mu_x, sig_x) * densite_normale(y, mu_y, sig_y)

# Test en (0,0) pour deux normales centrees reduites
d = densite_jointe_independante(0, 0, 0, 1, 0, 1)
theorique = (1 / math.sqrt(2 * math.pi))**2 # 1 / (2*pi)
print(f"Densite jointe en (0,0) : {d}")
assert abs(d - 1/(2*math.pi)) < 1e-5, "Le calcul de densite jointe est faux"
\end{lstlisting}


\end{document}
